\scanpage{172}
\uline{Corollaire 2}\quad Il y a une application unique
\[
 \widetilde C:Z^*(B/A)\mathrel{\text{---}}\Omega^*_{\widetilde B/\widetilde A}
\]
satisfaisant aux conditions suivantes

a) $\widetilde C(d\omega)=0$ pour \unsure{tte|tout} $\omega\in Z^*(B/A)$

b) $\widetilde C$ est un hom. de \deleted{\unsure{anneaux}}\inserted{\uline{$\Lambda$-algèbres}}
\[
 \left\{\begin{aligned}
 \widetilde C(\omega+\omega')&=C(\omega)+C(\omega')\\
 \widetilde C(\omega\omega')&=C(\omega)\,C(\omega')\\
 \widetilde C(\lambda\omega)&=\lambda C(\omega)\quad\text{si }\lambda\in\Lambda.\\
 C(1)&=1
 \end{aligned}\right.
\]

c) $\widetilde C(\deleted{\unsure{\xi}}x^{p-1}d_{B/A}x)=d_{\widetilde B/\widetilde A}x\otimes_{\widetilde B}1_\Lambda$

L'existence résulte de la définition. L'unicité résulte du corollaire 1.

\uline{Corollaire 3}\quad Si $x\in B^*$, on a
\[
 \text{d) }\ \widetilde C\left(\frac{dx}{x}\right)=\frac{dx}{x}\otimes_B1_\Lambda
 =\frac{d\widetilde x}{\widetilde x}\otimes_{\widetilde B}1_\Lambda
\]
En effet, on a $\frac{dx}{x}=x^{-p}(x^{p-1}dx)$, et comme $x^{-p}\in\Lambda$, on trouve
\[
 \begin{aligned}
 \widetilde C\left(\frac{dx}{x}\right)
 &=x^{-p}\unsure{\overline C|\widetilde C}(x^{p-1}dx)=\deleted{\unsure{x\,d\widetilde x}}F_A(x^{-1})C(x^{p-1}dx)\\
 &=(\widetilde x^{-1}d\widetilde x)\otimes_{\widetilde B}\Lambda
 \end{aligned}
\]

\uline{Cas particuliers}

a) Supposons que $\operatorname{Im}(A\mathrel{\text{--}} B)\deleted{\unsure{\supset}}\subset B^p$, \deleted{\unsure{i.e. $A=\operatorname{Im}(A\to B)$, Alors|i.e. $A\to\operatorname{Im}(A\to B)$. Alors}} p.ex $A$ parfait. Alors $\Omega^*_{B/A}\simeq\Omega^\alpha_{B/\mathbb F_p}\overset{\text{défn}}{=}\Omega^\alpha_B$, et on peut pour étudier \inserted{$H^*_{\mathrm{DR}}(B/A)$} supposer que $A=\mathbb F_p$. Mais, dans ce dernier cas, on aura $A\xrightarrow{\sim}\widetilde A$, donc $B\simeq\widetilde B$, et si $B\xrightarrow{F_B}B$ est injectif i.e. $\widetilde B\xrightarrow{\sim}\Lambda$ i.e. $B$ \uline{réduit}, alors le th. 1 revient à $\Omega^*_{B\deleted{\unsure{/A}}}\xrightarrow[\sim]{\varphi}H^*_{\mathrm{DR}}(B)$. Donc l'opération de Cartier peut être considérée alors comme
\[
 \overset{\deleted{\sim}}{C}:\deleted{\unsure{B}}Z^*(B)\longrightarrow\Omega^*_B
\]
On fera attention que la \inserted{dernière} formule s'écrit ici
\[
 \text{b')}\qquad \overset{\deleted{\sim}}{C}(f^p\omega)=fC(\omega)
\]
et c) et d) deviennent
% archive sheet number: 171

